% Generated draft: requires mathematical and source review.
\documentclass[11pt]{article}
\usepackage[margin=25mm]{geometry}
\usepackage{fontspec}
\setmainfont{Latin Modern Roman}
\usepackage{amsmath,amssymb,mathtools,unicode-math}
\setmathfont{Latin Modern Math}
\usepackage{enumitem,needspace,xurl,hyperref}
\setlength{\parindent}{0pt}
\setlength{\parskip}{5pt}
\setlength{\emergencystretch}{3em}
\newcommand{\N}{\mathbb{N}}
\newcommand{\Z}{\mathbb{Z}}
\newcommand{\Q}{\mathbb{Q}}
\newcommand{\R}{\mathbb{R}}
\newcommand{\C}{\mathbb{C}}
\newcommand{\F}{\mathbb{F}}
\newcommand{\A}{\mathbb{A}}
\newcommand{\PP}{\mathbb{P}}
\newcommand{\E}{\mathbb{E}}
\newcommand{\eps}{\varepsilon}
\newcommand{\dd}{\,\mathrm{d}}
\begin{document}
\section*{UnsolvedMath ID: 701178}
\textbf{Source problem code:} catalogue entry
\subsection{Definitions and mathematical statement}
Let \(\Gamma\) be a directed graph with infinitely many vertices. Write \(u\to v\) when there is an arc from \(u\) to \(v\). Let \(\overline{\Gamma}\) be the underlying undirected graph: distinct vertices are adjacent whenever there is an arc between them in either direction. Assume that \(\overline{\Gamma}\) is connected and has finite valency, and that some group \(G\leq\operatorname{Aut}(\Gamma)\) acts transitively on vertices, meaning that for every pair of vertices \(u,v\) there exists \(g\in G\) with \(g(u)=v\). Write \(d_{\overline{\Gamma}}(u,v)\) for the minimum length of an undirected path joining \(u\) and \(v\). Vertex transitivity implies that all vertices of \(\overline{\Gamma}\) have the same finite number \(d\) of distinct neighbours.

A directed path of length \(\ell\) is taken here to be a sequence of pairwise distinct vertices \(v_0,\ldots,v_\ell\) satisfying \(v_{i-1}\to v_i\) for \(1\leq i\leq\ell\). The source does not explicitly specify conventions concerning repeated vertices, loops, or parallel arcs; this draft uses the vertex-simple path convention and counts distinct neighbours for valency.

\textbf{Question (V. I. Trofimov, 2002).} For every directed graph \(\Gamma\) satisfying these hypotheses, does there exist \(k\in\Z_{>0}\), possibly depending on \(\Gamma\), such that
\[
\forall n\in\Z_{>0}\quad
\exists\ell\in\Z_{>0}\quad
\exists\text{ a directed path }(v_0,\ldots,v_\ell)\text{ in }\Gamma:
\qquad
\ell\leq kn
\quad\text{and}\quad
 d_{\overline{\Gamma}}(v_0,v_\ell)\geq n\,?
\]
The choice of \(k\) must precede the choice of \(n\); the path may depend on \(n\).

\textbf{Supplied background qualification.} The source's \emph{Comment of 2005} reports a quadratic bound: there exists a positive integer \(k'=k'(d)\), depending only on the valency \(d\) of \(\overline{\Gamma}\), such that for every \(n\in\Z_{>0}\) there is a directed path \((v_0,\ldots,v_\ell)\) in \(\Gamma\) with
\[
\ell\leq k'(d)n^2,
\qquad d_{\overline{\Gamma}}(v_0,v_\ell)\geq n.
\]
The dependence only on \(d\) means that the same coefficient may be used for all graphs satisfying the hypotheses with that valency. This is a result reported by the supplied source and attributed there to the 2006 article listed below; it is not independently verified here.
\subsection{Short English statement}
Does every infinite vertex-transitive directed graph with connected, finite-valency underlying undirected graph admit directed paths whose lengths are bounded linearly in the required undirected distance between their endpoints? The linear constant may depend on the graph. A supplied comment reports a quadratic bound with a constant depending only on valency.
\subsection{Sources}
\begin{itemize}
\item Local UnsolvedMath registry: supplied complete record, permanent ID \(701178\), exact source problem code \texttt{catalogue entry}, titled \emph{Kourovka Notebook Problem 15.90}. The record labels the problem \emph{open}; this is supplied registry status, not an independently verified literature assessment. The supplied record was updated on \texttt{2026-10-16T14:00:00Z}.
\item \emph{The Kourovka Notebook (Unsolved Problems in Group Theory)}, Problem 15.90, proposed by V. I. Trofimov in the 15th issue (2002). The supplied background also identifies the 21st issue edition (2026), \texttt{arXiv:1401.0300v48}. The collection description names E. I. Khukhro and V. D. Mazurov as editors.
\item The supplied \emph{Comment of 2005} attributes the quadratic bound to V. I. Trofimov, \emph{Europ. J. Combinatorics} \textbf{27}, no. 5 (2006), 690--700. No article title or article identifier is supplied.
\end{itemize}
\end{document}
